% GENERATED FILE. Edit canonical lessons, then run npm run generate:books.
% canonical-source-hash: fnv1a64:70fa88b73d4db74d
% canonical-lessons: TE-C214-bastandu, TE-C214-vimanasrayam, TE-C214-kudali, TE-C214-rajadhani, TE-C214-jilla

\chapter{Bus stand, Airport, Junction, Capital city, District}
\label{ch:te-bus-stand-airport-junction-capital-city-district}
\hlchaptermodality{214}

\begin{chapteropening}
\textbf{By the end of this chapter:} \emph{I can say a bus stand, an airport, a junction, a capital city and a district.}

Complete the last lesson of chapter 214: I can say a bus stand, an airport, a junction, a capital city and a district.
\end{chapteropening}

\section[\texorpdfstring{\te{బస్టాండు} (basṭāṇḍu)}{బస్టాండు (basṭāṇḍu)}]{\te{బస్టాండు} (basṭāṇḍu) — a bus stand}

\label{lesson:TE-C214-bastandu}



\begin{quote}
Before the new one: say the Telugu for a second, then the Telugu for the day before yesterday.
\end{quote}

\subsection*{You'll want to know: \te{బస్టాండు}}
\textbf{\te{బస్టాండు}} — \emph{basṭāṇḍu} — \textquotedblleft{}a bus stand\textquotedblright{}.

\begin{grammarlens}[title={What you've built}]
One more word for this chapter.
\end{grammarlens}

\subsection*{Your turn}
\begin{itemize}
\raggedright
  \item \emph{Say it:} \emph{basṭāṇḍu}
  \item \emph{Say it:} \emph{basṭāṇḍu}, once more
  \item \emph{Say it:} say \emph{basṭāṇḍu}
  \item \emph{Recall:} say the Telugu for a second, then the Telugu for the day before yesterday, then say \emph{basṭāṇḍu} again
\end{itemize}

\begin{tcolorbox}[breakable,colback=teal!4,colframe=teal!35!black,title={Before you move on}]
What does \te{బస్టాండు} mean? (\textquotedblleft{}A bus stand\textquotedblright{}.) Say it once more.
\end{tcolorbox}
\section[\texorpdfstring{\te{విమానాశ్రయం} (vimānāśrayaṁ)}{విమానాశ్రయం (vimānāśrayaṁ)}]{\te{విమానాశ్రయం} (vimānāśrayaṁ) — an airport}

\label{lesson:TE-C214-vimanasrayam}



\begin{quote}
Before the new one: say the Telugu for the day before yesterday, then the Telugu for a bus stand.
\end{quote}

\subsection*{You'll want to know: \te{విమానాశ్రయం}}
\textbf{\te{విమానాశ్రయం}} — \emph{vimānāśrayaṁ} — \textquotedblleft{}an airport\textquotedblright{}.

\begin{grammarlens}[title={What you've built}]
One more word for this chapter.
\end{grammarlens}

\subsection*{Your turn}
\begin{itemize}
\raggedright
  \item \emph{Say it:} \emph{vimānāśrayaṁ}
  \item \emph{Say it:} \emph{vimānāśrayaṁ}, once more
  \item \emph{Say it:} say \emph{vimānāśrayaṁ}
  \item \emph{Recall:} say the Telugu for the day before yesterday, then the Telugu for a bus stand, then say \emph{vimānāśrayaṁ} again
\end{itemize}

\begin{tcolorbox}[breakable,colback=teal!4,colframe=teal!35!black,title={Before you move on}]
What does \te{విమానాశ్రయం} mean? (\textquotedblleft{}An airport\textquotedblright{}.) Say it once more.
\end{tcolorbox}
\section[\texorpdfstring{\te{కూడలి} (kūḍali)}{కూడలి (kūḍali)}]{\te{కూడలి} (kūḍali) — a junction, a crossroads}

\label{lesson:TE-C214-kudali}



\begin{quote}
Before the new one: say the Telugu for a bus stand, then the Telugu for an airport.
\end{quote}

\subsection*{You'll want to know: \te{కూడలి}}
\textbf{\te{కూడలి}} — \emph{kūḍali} — \textquotedblleft{}a junction, a crossroads\textquotedblright{}.

\begin{grammarlens}[title={What you've built}]
One more word for this chapter.
\end{grammarlens}

\subsection*{Your turn}
\begin{itemize}
\raggedright
  \item \emph{Say it:} \emph{kūḍali}
  \item \emph{Say it:} \emph{kūḍali}, once more
  \item \emph{Say it:} say \emph{kūḍali}
  \item \emph{Recall:} say the Telugu for a bus stand, then the Telugu for an airport, then say \emph{kūḍali} again
\end{itemize}

\begin{tcolorbox}[breakable,colback=teal!4,colframe=teal!35!black,title={Before you move on}]
What does \te{కూడలి} mean? (\textquotedblleft{}A junction, a crossroads\textquotedblright{}.) Say it once more.
\end{tcolorbox}
\section[\texorpdfstring{\te{రాజధాని} (rājadhāni)}{రాజధాని (rājadhāni)}]{\te{రాజధాని} (rājadhāni) — a capital city}

\label{lesson:TE-C214-rajadhani}



\begin{quote}
Before the new one: say the Telugu for an airport, then the Telugu for a junction.
\end{quote}

\subsection*{You'll want to know: \te{రాజధాని}}
\textbf{\te{రాజధాని}} — \emph{rājadhāni} — \textquotedblleft{}a capital city\textquotedblright{}.

\begin{grammarlens}[title={What you've built}]
One more word for this chapter.
\end{grammarlens}

\subsection*{Your turn}
\begin{itemize}
\raggedright
  \item \emph{Say it:} \emph{rājadhāni}
  \item \emph{Say it:} \emph{rājadhāni}, once more
  \item \emph{Say it:} say \emph{rājadhāni}
  \item \emph{Recall:} say the Telugu for an airport, then the Telugu for a junction, then say \emph{rājadhāni} again
\end{itemize}

\begin{tcolorbox}[breakable,colback=teal!4,colframe=teal!35!black,title={Before you move on}]
What does \te{రాజధాని} mean? (\textquotedblleft{}A capital city\textquotedblright{}.) Say it once more.
\end{tcolorbox}
\section[\texorpdfstring{\te{జిల్లా} (jillā)}{జిల్లా (jillā)}]{\te{జిల్లా} (jillā) — a district}

\label{lesson:TE-C214-jilla}



\begin{quote}
Before the new one: say the Telugu for a junction, then the Telugu for a capital city.
\end{quote}

\subsection*{You'll want to know: \te{జిల్లా}}
\textbf{\te{జిల్లా}} — \emph{jillā} — \textquotedblleft{}a district\textquotedblright{}.

\begin{grammarlens}[title={What you've built}]
One more word for this chapter.
\end{grammarlens}

\subsection*{Your turn}
\begin{itemize}
\raggedright
  \item \emph{Say it:} \emph{jillā}
  \item \emph{Say it:} \emph{jillā}, once more
  \item \emph{Say it:} say \emph{jillā}
  \item \emph{Recall:} say the Telugu for a junction, then the Telugu for a capital city, then say \emph{jillā} again
\end{itemize}

\begin{tcolorbox}[breakable,colback=teal!4,colframe=teal!35!black,title={Before you move on}]
What does \te{జిల్లా} mean? (\textquotedblleft{}A district\textquotedblright{}.) Say it once more.
\end{tcolorbox}
